\begin{abstract}
Approximately 1.4 million active-duty military personnel and roughly 600,000 overseas civilians and dependents are counted in the decennial Census at their home-of-record state rather than their current physical location. This counting convention---required by 13 U.S.C. \S\,141 and implemented through the Census Bureau's Federal Employees Overseas Rule---concentrates military-affiliated population credits in a small number of high-military-presence states and, within those states, in a small number of congressional districts that contain major installations. We analyze the redistricting effects of military population counting in the three states most affected: Virginia, North Carolina, and Texas. In North Carolina, the Fayetteville metropolitan area surrounding Fort Liberty (formerly Fort Bragg) concentrates approximately 100,000 active-duty and dependent personnel whose home-of-record credits flow disproportionately to the surrounding congressional district. We estimate that this concentration inflates the military-area district by approximately 6--8\% relative to the civilian-only counterfactual, requiring boundary compression that benefits adjacent districts. BISECT uses Census total population, which includes all domestically stationed military personnel but excludes overseas personnel counted at home-of-record state (not tract). We analyze the partisan direction of military population counting: the active military leans Republican by a margin of approximately 12--18 percentage points in recent surveys, creating a mild systematic R-advantage in high-military districts. The Uniformed and Overseas Citizens Absentee Voting Act (UOCAVA) governs voting rights for this population. We identify three specific policy choices in BISECT's pipeline that interact with military population counting and propose a \texttt{--balance-metric civilian\_only} sensitivity flag for research use.
\end{abstract}
